\documentclass[9pt,english,]{beamer}
\usepackage{lmodern}
\usepackage{amssymb,amsmath}
\usepackage{ifxetex,ifluatex}

\ifnum 0\ifxetex 1\fi\ifluatex 1\fi=0 % if pdftex
  \usepackage[T1]{fontenc}
  \usepackage[utf8]{inputenc}
\else % if luatex or xelatex
  \ifxetex
    \usepackage{mathspec}
  \else
    \usepackage{fontspec}
  \fi
  \defaultfontfeatures{Ligatures=TeX,Scale=MatchLowercase}
\fi
% use upquote if available, for straight quotes in verbatim environments
\IfFileExists{upquote.sty}{\usepackage{upquote}}{}
% use microtype if available
\IfFileExists{microtype.sty}{%
\usepackage{microtype}
\UseMicrotypeSet[protrusion]{basicmath} % disable protrusion for tt fonts
}{}
\usepackage{hyperref}
\PassOptionsToPackage{usenames,dvipsnames}{color} % color is loaded by hyperref
\hypersetup{unicode=true,
            pdftitle={Elements of Matrix Algebra},
            pdfauthor={Yannick Hoga},
            colorlinks=true,
            linkcolor=dueblue,
            citecolor=dueblue,
            urlcolor=dueblue,
            breaklinks=true}
\urlstyle{same}  % don't use monospace font for urls
\ifnum 0\ifxetex 1\fi\ifluatex 1\fi=0 % if pdftex
  \usepackage[shorthands=off,main=english]{babel}
\else
  \usepackage{polyglossia}
  \setmainlanguage[]{}
\fi


\usepackage{longtable,booktabs}
\IfFileExists{parskip.sty}{%
\usepackage{parskip}
}{% else
\setlength{\parindent}{0pt}
\setlength{\parskip}{3em plus 1pt minus 1pt}%{6pt plus 2pt minus 1pt}
}
\setlength{\emergencystretch}{3em}  % prevent overfull lines

\setcounter{secnumdepth}{0}
% Redefines (sub)paragraphs to behave more like sections
\ifx\paragraph\undefined\else
\let\oldparagraph\paragraph
\renewcommand{\paragraph}[1]{\oldparagraph{#1}\mbox{}}
\fi
\ifx\subparagraph\undefined\else
\let\oldsubparagraph\subparagraph
\renewcommand{\subparagraph}[1]{\oldsubparagraph{#1}\mbox{}}
\fi

%%% Use protect on footnotes to avoid problems with footnotes in titles
\let\rmarkdownfootnote\footnote%
\def\footnote{\protect\rmarkdownfootnote}


  \title{Elements of Matrix Algebra}
    \author{Yannick Hoga}
  
  \providecommand{\institute}[1]{}
  \expandafter\ifstrequal\expandafter{default}{default}{%
  \institute{\expandafter\ifstrequal\expandafter{en}{de}{%
  Universität Duisburg-Essen}{%
  University of Duisburg-Essen}}}{%
  
  \institute{default}
}


    \date{Winter 2024}


% up until here https://github.com/rstudio/rmarkdown/blob/master/inst/rmd/latex/default-1.17.0.2.tex
%%%%%%%%%%%%%%%%%%%%%%%%%%%%%%%%%%%%%%%%%%%%%%%%%%%%%%%%%%%%%%%%%%%%%%%%%%%%%%%
%%% Additional packages
\makeatletter
\def\input@path{{"C:/Users/Yannick.Hoga/AppData/Local/Programs/R/R-4.3.3/library/runidue/rmarkdown/templates/lectureslides/resources/"}}  % make them visible to tex to overcome warnings
\makeatother

\input{"C:/Users/Yannick.Hoga/AppData/Local/Programs/R/R-4.3.3/library/runidue/rmarkdown/templates/lectureslides/resources/ee.sty"}
\usepackage{BeamerColor}
\usepackage{hypernat}
\usepackage{mathtools}
\usepackage[export]{adjustbox}
\usepackage{makecell}
\usepackage{diagbox}
\usepackage{bm, bbm}
\usepackage{subfig}
\usepackage{forest}
\usepackage[labelfont={footnotesize,color=dueblue},font=footnotesize]{caption}
\usepackage{soul}
\usepackage{textcomp}
\usepackage{zref}
\usepackage{zref-user}
\usepackage{zref-xr}
\usepackage{empheq}
\usepackage{tikz}
\usetikzlibrary{arrows, decorations.markings}
\tikzstyle{arrow} = [ultra thick, - latex, shorten >= 2pt, shorten <= 2pt] % default arrow
\tikzstyle{box} = [rectangle,rounded corners,fill=dueblue, text=white, align=center, node distance=1.25cm] % default node

% \usepackage{etoolbox}

\makeatletter
\let\@@magyar@captionfix\relax
\makeatother

%%% ZREF

%%% New colors
\definecolor{dueblue}{RGB}{0,76,147}
\definecolor{duelightblue}{RGB}{223,228,242}
\definecolor{duebeige}{RGB}{239,228,191}
\definecolor{darkgrey}{RGB}{40, 40, 40}
\definecolor{salmon}{RGB}{239, 148, 108}
\definecolor{lime}{RGB}{219, 254, 135}
\definecolor{sandy}{RGB}{255, 227, 129}
\definecolor{darklime}{RGB}{153, 194, 77}


%%% Additional commands
\newcommand{\npt}{\\[1ex]\pause}
\newcommand{\nbs}{\protect{~}}
\newcommand{\qedb}{\hfill \textsquare} % to be used outside proof environments
\newcommand{\pausealt}{ }
\newcommand{\blue}[1]{{\color{dueblue}{#1}\color{black}}}
\newcommand{\hil}[1]{\textbf{\color{dueblue}{#1}\color{black}}}
\newcommand{\hilm}[1]{\mathbf{\color{dueblue}{#1}\color{black}}}
\newcommand{\mps}{\par\medbreak}
\newcommand{\sd}{\operatorname{sd}}
\newcommand{\SD}{\operatorname{SD}}

% Kahoot Block
\newenvironment<>{kahoot}[1]{%
  \begin{actionenv}#2%
      \def\insertblocktitle{#1}%
      \par%
      \mode<presentation>{%
      \setbeamercolor{block title}{fg=black,bg=yellow}
       \setbeamercolor{block body}{fg=black,bg=yellow!20}
       \setbeamercolor{itemize item}{fg=orange!20!black}
       \setbeamertemplate{itemize item}[triangle]
     }%
      \usebeamertemplate{block begin}
      \begin{center}
\includegraphics[width=0.25\textwidth]{"C:/Users/Yannick.Hoga/AppData/Local/Programs/R/R-4.3.3/library/runidue/rmarkdown/templates/lectureslides/resources/kahootlogo2"}
\end{center}}
    {\par\usebeamertemplate{block end}\end{actionenv}}
\newcommand{\kahootblock}[1]{\begin{kahoot}{#1}\end{kahoot}}

\makeatletter
\newcommand*\fsize{\dimexpr\f@size pt\relax} % returns current font size
\makeatother
\newcommand{\R}{\includegraphics[height=.85\fsize, clip=T]{"C:/Users/Yannick.Hoga/AppData/Local/Programs/R/R-4.3.3/library/runidue/rmarkdown/templates/lectureslides/resources/Rlogo"}\ }
\newcommand{\btVFill}{\vskip0pt plus 1filll}
\newcommand{\source}[1]{\btVFill\scriptsize \textit{\expandafter\ifstrequal\expandafter{en}{de}{Quelle: }{Source: }}#1}



% BLOCK ENVIRONMENTS (THEOREMS)
\expandafter\ifstrequal\expandafter{box}{blank}{
}{
\setbeamerfont{block title}{size=\normalsize}
\setbeamercolor{block title}{bg=duelightblue,fg=black}
\setbeamercolor{block body}{bg=duelightblue!37,fg=black}
}%
\setbeamertemplate{blocks}[rounded][shadow=false]

% Headers for theorems
\expandafter\ifstrequal\expandafter{en}{de}{%
  \newcommand{\thmncont}{Fortsetzung}
  \newtheorem{thmn}{Theorem}[section]
  \newtheorem{lem}[thmn]{Lemma}
  \newtheorem{cor}[thmn]{Corollary}
  \newtheorem{defn}[thmn]{Definition}
  \newtheorem{ass}[thmn]{Annahme}
  \newtheorem{rem}[thmn]{Anmerkung}
  \newtheorem{prop}[thmn]{Proposition}
  \theoremstyle{definition}
  \newtheorem{xmpl}[thmn]{Beispiel}
  \newtheorem{que}[thmn]{Frage}
  \newtheorem{exe}[thmn]{Aufgabe}
  \numberwithin{equation}{section}}{%%%
  \newcommand{\thmncont}{Continued}
  \newtheorem{thmn}{Theorem}[section]
  \newtheorem{lem}[thmn]{Lemma}
  \newtheorem{cor}[thmn]{Corollary}
  \newtheorem{defn}[thmn]{Definition}
  \newtheorem{ass}[thmn]{Assumption}
  \newtheorem{rem}[thmn]{Remark}
  \newtheorem{prop}[thmn]{Proposition}
  \theoremstyle{definition}
  \newtheorem{xmpl}[thmn]{Example}
  \newtheorem{que}[thmn]{Question}
  \newtheorem{exe}[thmn]{Exercise}
  \numberwithin{equation}{section}
}%



\setbeamertemplate{theorems}[numbered]
\setbeamerfont{block title}{size=\normalsize}
\setbeamertemplate{theorem begin}
{%
  \expandafter\ifstrequal\expandafter{\inserttheoremaddition}{*}{\addtocounter{thmn}{-1}}{}
  \normalfont
  \begin{\inserttheoremblockenv}{%
    \structure{%
      \textcolor{darkgrey}{\textbf{\inserttheoremname \inserttheoremnumber}:
      \ifx\inserttheoremaddition\empty\else\expandafter\ifstrequal\expandafter{\inserttheoremaddition}{*}{\thmncont\inserttheorempunctuation}{\inserttheoremaddition\inserttheorempunctuation}\fi}%
    }%
  }%
}
\setbeamertemplate{theorem end}{\leavevmode\end{\inserttheoremblockenv}}

\preto{\xmpl}{
    \setbeamercolor{block title}{fg=black,bg=sandy!70!white}
    \setbeamercolor{block body}{fg=black, bg=sandy!20!white}
}

\preto{\exe}{
    \setbeamercolor{block title}{fg=black,bg=darklime!70!white}
    \setbeamercolor{block body}{fg=black, bg=darklime!20!white}
}




% Spacings
\AtBeginEnvironment{figure}{\setlength{\parskip}{-2ex}}
\AtBeginEnvironment{center}{\smallbreak}
\AtEndEnvironment{center}{\smallbreak}
\AtBeginEnvironment{itemize}{\smallskip \setlength{\itemsep}{0.3ex}\setlength{\parskip}{1ex}}
\AtEndEnvironment{itemize}{\smallskip}
\preto{\enumerate}{\medskip  \setlength{\itemsep}{0.75ex}\setlength{\parskip}{1ex}}
\AtEndEnvironment{enumerate}{\smallskip}



% \let\tempone\itemize
% \let\temptwo\enditemize
% \renewenvironment{itemize}{\tempone\setlength{\itemsep}{1ex}\setlength{\parskip}{1.2ex}}{\temptwo}

% \let\tempone\enumerate
% \let\temptwo\endenumerate
% \renewenvironment{enumerate}{\tempone\setlength{\itemsep}{1ex}\setlength{\parskip}{1.2ex}}{\temptwo}
% 




%%% LAYOUT
\usecolortheme[named=dueblue]{structure}
\beamertemplatenavigationsymbolsempty
\setbeamersize{text margin left=10pt, text margin right=10pt}


% Spacing before and after code chunks / output
\makeatletter
\preto{\@verbatim}{\topsep=-2pt \partopsep=2pt \parskip=-2ex}
\makeatother

% spacing before figure captions
\setlength{\abovecaptionskip}{0pt}
\setlength{\belowcaptionskip}{-1em plus 1em minus 0em}

% spacing aorund display equations
\setlength{\abovedisplayskip}{3pt}
%\setlength{\belowdisplayskip}{3pt}

 % 

% Frame Head
\makeatletter
\setbeamertemplate{frametitle}{%
\vspace{.1cm}
\begin{minipage}[t][1cm]{.84\textwidth}
    \ifbeamer@plainframe%
      \Large\insertframetitle%
    \else
      \Large\insertframetitle%
      \ifx\insertframesubtitle\empty
      \else
        \\ \small \insertframesubtitle
      \fi
    \fi
\end{minipage}\hfill


\ifstrequal{FALSE}{FALSE}{}{\begin{minipage}[t][1cm]{1.9cm}\includegraphics[width=2cm,keepaspectratio=true, valign=t]{FALSE}\end{minipage}\vspace{-.3cm}}


}
\makeatother

% Footer
\setbeamercolor{footlinecolor}{bg=dueblue,fg=white}


\ifstrequal{colored}{colored}{
\setbeamertemplate{footline}{
  % \hbox{
  \begin{beamercolorbox}[wd=\paperwidth,ht=0.4cm, leftskip=0cm, rightskip=0cm]{footlinecolor}
    \hfill \scriptsize{{\insertshortauthor : }
    \insertshorttitle---\insertsection \hfill \thesection-\insertframenumber\hspace*{.15cm}}\vspace{1mm}
  \end{beamercolorbox}
  % }
}
\addtobeamertemplate{footline}{\hypersetup{linkcolor=.}}{}
}{
\setbeamertemplate{footline}{\hspace*{.1cm}%
\hfill \scriptsize{{\insertshortauthor : }
\insertshorttitle---\insertsection \hfill \thesection-\insertframenumber\hspace*{.2cm}}\vspace{1mm}}
}




% TOC


\setbeamertemplate{section in toc shaded}[default][50]
\defbeamertemplate{section in toc}{customcircle}
{\leavevmode\leftskip=4ex%
  \llap{%
    \usebeamerfont*{section number projected}%
    \usebeamercolor{section number projected}%
    \begin{pgfpicture}{-1ex}{-.2ex}{1ex}{2ex}
      \color{bg}
      \pgfpathcircle{\pgfpoint{0pt}{.75ex}}{1.3ex}
      \pgfusepath{fill}
      \pgftext[base]{\color{fg}\inserttocsectionnumber}
    \end{pgfpicture}\kern1.25ex%
  }%
  \hspace{0.8ex}\inserttocsection\par}

\setbeamertemplate{section in toc}[customcircle]

% \usepackage{xpatch}
% \xpatchcmd{\itemize}
%   {\def\makelabel}
%   {\setlength{\itemsep}{1ex}\def\makelabel}
%   {}
%   {}
% \makeatletter
% \xpatchcmd{\beamer@enum@}
%   {\def\makelabel}
%   {\setlength{\itemsep}{1ex}\def\makelabel}
%   {}
%   {}
% \makeatother

% NEW SECTION


\AtBeginSection[] {
{
\setbeamertemplate{footline}{}
\begin{frame}{\expandafter\ifstrequal\expandafter{en}{de}{Überblick}{Outline}}
\vspace{0.5cm}
\tableofcontents[currentsection]
\end{frame}
\addtocounter{framenumber}{-\insertframenumber} % collect all begin section commands in one env.!
}
}

% ITEMIZE, ENUMERATE, DESCRIPTION
\defbeamertemplate*{description item}{defaultperiod}{\insertdescriptionitem.}
\setbeamersize{description width=15pt} % left margin for description

\setlength{\topsep}{1ex}
\setlength{\leftmargini}{10pt+.5\labelwidth}
\setlength{\leftmarginii}{15pt}
\setlength{\leftmarginiii}{15pt}


\setbeamertemplate{caption}[numbered]
\setbeamertemplate{caption label separator}{: }
\setbeamertemplate{itemize item}[circle]
\setbeamertemplate{itemize subitem}[triangle]
\setbeamertemplate{enumerate item}[default]
\setbeamertemplate{enumerate subitem}[default]
\setbeamertemplate{description item}[defaultperiod]
\setbeamertemplate{description subitem}[defaultperiod]


\providecommand{\tightlist}{}% \setlength{\itemsep}{0.75ex}\setlength{\parskip}{1.2ex}
% allow environments inside environments
\let\otherbegin\begin
\let\otherend\end
%%%%%%%%%%%%%%%%%%%%%%%%%%%%%%%%%%%%%%%%%%%%%%%%%%%%%%%%%%%%%%%%%%%%%%%%%%%%%%%
% from here on https://github.com/rstudio/rmarkdown/blob/master/inst/rmd/latex/default-1.17.0.2.tex



\begin{document}


      {\usebackgroundtemplate{\includegraphics[scale=0.25]{"C:/Users/Yannick.Hoga/AppData/Local/Programs/R/R-4.3.3/library/runidue/rmarkdown/templates/lectureslides/resources/duetower"}}
    \frame[plain]{\titlepage}}
  

% 
\begin{frame}{Introduction}
\protect\hypertarget{introduction}{}
The material covered here may be found in numerous appendices of
econometrics textbooks. We will only give definitions and results
relevant to econometrics (you will be asked to verify some in the
exercises). For more, see, e.g.,

\begin{itemize}
\tightlist
\item
  Rao and Rao (1998) \emph{Matrix Algebra and Its Applications to
  Statistics and Econometrics}. World Scientific, 1st ed.

  \begin{itemize}
  \tightlist
  \item
    Has the necessary background on linear algebra.
  \end{itemize}
\item
  Magnus and Neudecker (2019) \emph{Matrix Differential Calculus with
  Applications in Statistics and Econometrics}. Wiley, 3rd ed.

  \begin{itemize}
  \tightlist
  \item
    The title summarizes it aptly.
  \end{itemize}
\item
  Lütkepohl (1996) \emph{Handbook of Matrices}. Wiley, 3rd ed.

  \begin{itemize}
  \tightlist
  \item
    Has a lot more results for ready reference.
  \end{itemize}
\end{itemize}
\end{frame}

\begin{frame}{Definitions, Notation, Terminology}
\protect\hypertarget{definitions-notation-terminology}{}
\defn[Matrix]A \((m \times n)\) \hil{matrix} \(\mA\) is an array of
numbers: \[
\mA=\left(%
\begin{array}{ccc}
  a_{11} & \cdots & a_{1n} \\
  \vdots & \ddots & \vdots \\
  a_{m1} & \cdots & a_{mn}
\end{array}%
\right)
\] Alternative notation: \(\mA\;(m \times n)\),
\(\mA=[a_{ij}]\;(m\times n)\). \medskip

\(m\) and \(n\) are positive integers denoting the \hil{row dimension}
and the \hil{column dimension}, respectively, \((m \times n)\) is the
\hil{dimension} or order of \(\mA\).\medskip

For our purposes, the \(a_{ij},\;i = 1,\ldots, m,\;j = 1,\ldots, n\),
are elements of \(\mathbb{R}\).\medskip

As such, \(\mA\) can be viewed as a linear mapping
\(\mA:\mathbb{R}^n\to\mathbb{R}^m\). \enddefn
\end{frame}

\begin{frame}{}
\protect\hypertarget{section}{}
\defn[Block Matrix]\[
\mA= \left(%
\begin{array}{ccc}
  \mA_{11} & \cdots & \mA_{1n} \\
  \vdots & \ddots & \vdots \\
  \mA_{m1} & \cdots & \mA_{mn} \\
\end{array}%
\right)
\] Here, the \(\mA_{ij}\) are \((m_i \times n_j)\) submatrices of
\(\mA\). A matrix written in terms of submatrices rather than individual
elements is often called a \hil{partitioned matrix} or a
\hil{block matrix}. \enddefn
\end{frame}

\begin{frame}{}
\protect\hypertarget{section-1}{}
\defn[Special Matrices and Elements]A matrix \(\mA\;(m \times m)\)
having the same number of rows and columns is a \hil{square} or
\hil{quadratic matrix}.\medskip

The elements \(a_{ii}, \;i= 1,\ldots,m,\) constitute its
\hil{principal diagonal}.\medskip

If \(a_{ii}=1\) for \(i = 1,\ldots, m\), and \(a_{ij}= 0\) for
\(i\neq j\) is an (\(m \times m)\) \hil{identity matrix}, denoted
\(\mI_m\).\medskip

A \((1\times n)\) matrix \(\va\) is an \(n\)-dimensional
\hil{row vector} or \((1\times n)\) vector.\medskip

An \((m \times 1)\) matrix \(\vb\) is an \(m\)-dimensional
\hil{column vector}. \enddefn
\end{frame}

\begin{frame}{}
\protect\hypertarget{section-2}{}
\defn[Trace and Transpose]The \hil{trace} of \(\mA\) is defined as \[
  \tr(\mA) = \sum_{i=1}^ma_{ii}
\] The \hil{transpose} of \(\mA\;(m \times n)\) is defined as \[
\mA^\top=\left(%
\begin{array}{ccc}
  a_{11} & \cdots & a_{1n} \\
  \vdots & \ddots & \vdots \\
  a_{m1} & \cdots & a_{mn}
\end{array}%
\right)^\top=\left(%
\begin{array}{ccc}
  a_{11} & \cdots & a_{m1} \\
  \vdots & \ddots & \vdots \\
  a_{1n} & \cdots & a_{mn}
\end{array}%
\right)
\] \(\mA^\top\) is \((n \times m)\). That is, the rows of \(\mA\) are
the columns of \(\mA^\top\). \enddefn
\end{frame}

\begin{frame}{}
\protect\hypertarget{section-3}{}
\defn[Determinant]The \hil{determinant} of \(\mA\;(m\times m)\)
\begin{eqnarray*}
 \det(\mA) =|\mA|&=& a_{i1} \textrm{cof}(a_{i1})+\ldots+
a_{im}\textrm{cof}(a_{im}) \\&=&
a_{1j}\textrm{cof}(a_{1j})+\ldots+ a_{mj}\textrm{cof}(a_{mj}),
\end{eqnarray*} where
\(\textrm{cof}(a_{ij})=(-1)^{i+j}\textrm{minor}(a_{ij})\) and \[
\textrm{minor}(a_{ij})=\left|%
\begin{array}{cccccc}
  a_{1,1} & \cdots & a_{1,j-1} &  a_{1,j+1} & \cdots &  a_{1,m}\\
  \vdots &  & \vdots & \vdots &  & \vdots \\
  a_{i-1,1} & \cdots & a_{i-1,j-1} & a_{i-1,j+1} & \cdots & a_{i-1,m} \\
   a_{i+1,1} & \cdots & a_{i+1,j-1} & a_{i+1,j+1} & \cdots & a_{i+1,m} \\
  \vdots &  & \vdots & \vdots &  & \vdots \\
  a_{m,1} & \cdots & a_{m,j-1} & a_{m,j+1} & \cdots & a_{m,m} \\
\end{array}%
\right|
\] \enddefn
\end{frame}

\begin{frame}{}
\protect\hypertarget{section-4}{}
\defn[Inverses]For \(\mA \;(m \times m)\) with \(\det(\mA)\neq 0\), the
\hil{inverse} of \(\mA\) is the unique \((m \times m)\) matrix
\(\mA^{-1}\) satisfying \(\mA\mA^{-1} = \mA^{-1}\mA = \mI_m\).\medskip

An \((n \times m)\) matrix \(\mA^-\) is a \hil{generalized inverse} of
the \((m \times n)\) matrix \(\mA\) if it satisfies
\(\mA\mA^-\mA = \mA\). \(\mA^-\) is generally not unique.\medskip

The \((n \times m)\) matrix \(\mA^+\) is the
\hil{Moore--Penrose (generalized) inverse} of the \((m \times n)\)
matrix \(\mA\) if it satisfies

\begin{itemize} 
  \item $\mA\mA^+\mA = \mA$,
  \item $\mA^+\mA\mA^+ = \mA^+$, 
  \item $(\mA\mA^+)^\top = \mA\mA^+$, 
  \item $(\mA^+\mA)^\top = \mA^+\mA$ .
\end{itemize}

\(\mA^+\) exists and is unique. \enddefn
\end{frame}

\begin{frame}{}
\protect\hypertarget{section-5}{}
\defn[Certain Matrix Products]\hil{Power} of a matrix
\(\mA\;(m \times m)\): \[ \mA^i=\begin{cases} \prod_{j=1}^i \mA =
\underbrace{\mA\times
\cdots \times \mA}_{i \textrm{ times}} & \text{for positive integers, }i\in\mathbb{Z}^+ \\[.5cm] \mI_m &\textrm{for } i=0 \\[.5cm]
(\prod_{j=1}^{-i} \mA)^{-1} & \text{for }i\in\mathbb{Z}^-,
\text{if } \det(\mA)\neq 0
\end{cases}
\] The \((m \times m)\) matrix \(\mA^{1/2}\) is a \hil{square root} of
the \((m \times m)\) matrix \(\mA\) if \(\mA^{1/2}\mA^{1/2}=\mA\).
\enddefn
\end{frame}

\begin{frame}{}
\protect\hypertarget{section-6}{}
\defn[Vectorization]The \hil{vectorization} of \(\mA\;(m \times n)\) is
\[
\vec(\mA)=\left(%
\begin{array}{c}
  a_{11} \\
  \vdots \\
  a_{m1} \\
  a_{12} \\
  \vdots \\
  a_{m2} \\
  \vdots \\
  a_{1n} \\
  \vdots \\
  a_{mn} \\
\end{array}%
\right)\;(mn \times 1) \] That is, vec stacks the columns of \(\mA\) in
a column vector. \enddefn
\end{frame}

\begin{frame}{}
\protect\hypertarget{section-7}{}
\defn[Some Matrix Properties]\zlabel{matrixproperties}
\begin{enumerate}
\item\zlabel{it:9.1} An $(m \times m)$ matrix with $\mA = \mA^\top$ is \hil{symmetric}.

\item\zlabel{it:9.2} An $(m \times m)$ matrix $\mA$ is \hil{idempotent} if $\mA^2 = \mA$.

\item\zlabel{it:9.3} An $(m \times m)$ matrix $\mA$ is \hil{nonsingular} or \hil{invertible} or \hil{regular} if $\det(\mA) \neq 0$ and thus $\mA^{-1}$ exists.

\item\zlabel{it:9.4} An $(m \times m)$ matrix $\mA$ is \hil{orthogonal} if $\mA$ is nonsingular and $\mA^\top = \mA^{-1}$.

\end{enumerate}
\enddefn
\end{frame}

\begin{frame}{Concepts Related to Individual Matrices}
\protect\hypertarget{concepts-related-to-individual-matrices}{}
\defn[Definiteness]\zlabel{Definiteness}

A symmetric \((m \times m)\) matrix \(\mA\) is

\begin{itemize}

\item \hil{positive definite} if $\vx^\top \mA\vx > 0$ for all $(m\times 1)$ vectors $\vx \neq \vzero$ 

\item \hil{positive semidefinite} if $\vx^\top \mA\vx \geqslant 0$ for all $(m \times 1)$ vectors $\vx$

\item \hil{negative definite} if $\vx^\top \mA\vx < 0$ for all $(m\times 1)$ vectors $\vx \neq \vzero$

\item \hil{negative semidefinite} if $\vx^\top \mA\vx \leqslant 0$ for all $(m \times 1)$ vectors $\vx$

\item \hil{indefinite} if there exist $(m \times 1)$ vectors $\vx$ and $\vy$ such that $\vx^\top \mA\vx > 0$ and $\vy^\top \mA\vy < 0 $.

\end{itemize}
\enddefn
\end{frame}

\begin{frame}{}
\protect\hypertarget{section-8}{}
\defn[Triangular Matrices]An \((m \times m)\) matrix \[
\mA=\left(%
\begin{array}{cccc}
  a_{11} & 0 & \cdots & 0 \\
  \vdots & \ddots & \ddots & \vdots \\
  \vdots &  & \ddots & 0 \\
  a_{m1} & \cdots & \cdots & a_{mm} \\
\end{array}%
\right)
\] is \hil{lower triangular}.\medskip

A matrix \(\mA\;(m\times m)\) that is both upper and lower triangular is
a \hil{diagonal matrix}, denoted \(\mA = \diag(a_{11},\ldots,a_{mm})\).
\enddefn
\end{frame}

\begin{frame}{Further Concepts}
\protect\hypertarget{further-concepts}{}
\defn\zlabel{linindep}

The \(m\)-dimensional row or column vectors \(\vx_1,\ldots,\vx_k\) are
\hil{linearly independent} if, for scalars \(c_1,\ldots,c_k\),
\(c_1\vx_1+\ldots+c_k \vx_k = \vzero\) implies \(c_1,\ldots,c_k = 0\).
\enddefn

\defn

The \hil{rank} of \(\mA\;(m\times n)\), denoted \(\rk(\mA)\), is the
maximum number of linearly independent rows or columns of \(\mA\).
\enddefn

\defn\zlabel{quadrform}

Given a symmetric \((m \times m)\) matrix \(\mA\), the function
\(Q: \mathbb{R}^m\rightarrow\mathbb{R}\) defined by
\(Q(\vx) = \vx^\top \mA\vx\) is called a \hil{quadratic form}. \enddefn

\defn

The polynomial in \(\mA\) given by \(\det(\lambda\,\mI_m - \mA)\) is the
\hil{characteristic polynomial} of the \((m \times m)\) matrix \(\mA\).
\enddefn
\end{frame}

\begin{frame}{}
\protect\hypertarget{section-9}{}
\defn[Eigenvalues and Eigenvectors]The \hil{eigenvalues} of an
\((m \times m)\) matrix \(\mA\) are the roots of the characteristic
polynomial of \(\mA\).\medskip

(Note that eigenvalues can be complex depending on the
polynomial.)\medskip

An \((m \times 1)\) vector \(\vv \neq \vzero\) satisfying
\(\mA\vv = \lambda \vv\), where \(\lambda\) is an eigenvalue of \(\mA\),
is an \hil{eigenvector} corresponding to the eigenvalue \(\lambda\).
\enddefn
\end{frame}

\begin{frame}{Kronecker Product}
\protect\hypertarget{kronecker-product}{}
\defn

The \hil{Kronecker product} of \(\mA\;(m\times n)\) and
\(\mB\;(p\times q)\) is defined as the \((mp \times nq)\) matrix \[
  \mA\otimes \mB=\left(%
\begin{array}{ccc}
  a_{11}\mB & \cdots & a_{1n}\mB \\
  \vdots & \ddots & \vdots \\
  a_{m1}\mB & \cdots & a_{mn}\mB
\end{array}%
\right)
\] \enddefn
\end{frame}

\begin{frame}{Operations Relating Matrices}
\protect\hypertarget{operations-relating-matrices}{}
\defn
\begin{enumerate} 

\item \hil{Addition:} $\mA \;(m \times n),\;\mB\;(m \times n): \mA+\mB= [a_{ij}+b_{ij}]\;(m \times n)$ 

\item $\mA,\mB,\mC\;(m\times n) : (\mA \pm
\mB) \pm \mC = \mA \pm (\mB \pm \mC)$ 
  
\item \hil{Multiplication by a scalar:} $\mA\;(m \times n),\;c$ a scalar $c\mA = [c\cdot a_{ij}]\;(m \times
n)$

\item \hil{Matrix multiplication} or matrix product: $\mA \; (m \times
n),\; \mB\;(n \times p)$
\[
  \mA\mB_{ij}=\sum_{k=1}^na_{ik}b_{kj}\;\;\forall\;i=1,\ldots,m;\;j=1,\ldots,p\;(m \times p) 
\]
If the blocks are conformable, multiplication of partitioned matrices proceeds analogously.

\end{enumerate}
\enddefn
\end{frame}

\begin{frame}{Rules}
\protect\hypertarget{rules}{}
\prop\zlabel{rules1}
\begin{enumerate}
\item\zlabel{it:19.1} $\mA\;(m \times n)$ : $\mI_m\mA = \mA\mI_n=\mA$

\item\zlabel{it:19.2} $\mA\; (m\times n),\mB,\mC\;(n \times r)$ : $\mA(\mB \pm \mC) = \mA\mB \pm \mA\mC$

\item\zlabel{it:19.3} $\mA\;(m\times n),\mB,\mC \;(p\times q)$ : $\mA \otimes (\mB \pm \mC) = \mA \otimes \mB \pm \mA \otimes \mC$

\item\zlabel{it:19.4} $\mA,\mB\;(m\times n)$ : $(\mA\pm \mB)^\top = \mA^\top\pm \mB^\top$

\item\zlabel{it:19.5} $\mA, \mB\; (m \times m)$ : $\tr(\mA \pm \mB) = \tr(\mA)\pm \tr(\mB)$

\item\zlabel{it:19.6} $\mA, \mB\; (m \times n)$ : $\rk(\mA \pm \mB) \leqslant \rk(\mA) + \rk(\mB)$

\item\zlabel{it:19.7} $\mA\;(m\times n),\mB\;(n\times p)$ : $\rk(\mA\mB) \leqslant \min(\rk(\mA),\rk(\mB))$

\item\zlabel{it:19.8} $\mA \;(m \times m)$ : $\mA + \mA^\top$ is a symmetric $(m \times m)$ matrix

\end{enumerate}
\endprop
\end{frame}

\begin{frame}{Rules (2)}
\protect\hypertarget{rules-2}{}
\prop\zlabel{rules2}
\begin{enumerate}
\item\zlabel{it:20.1} $(\mA\mB)\mC = \mA(\mB\mC) = \mA\mB\mC$

\item\zlabel{it:20.2} $\mA\mB \neq \mB\mA$ in general

\item\zlabel{it:20.3}  $\mA\;(m \times n), \mB\;(p\times q),\mC\;(n\times r),\mD\;(q \times s) : (\mA \otimes \mB)(\mC \otimes \mD) = \mA\mC\otimes \mB\mD$.

\item\zlabel{it:20.4} $(\mA\mB)^\top = \mB^\top \mA^\top$
    
\item\zlabel{it:20.5} $(\mA\mB)^{-1}= \mB^{-1}\mA^{-1}$

\item\zlabel{it:20.6} $\mA\;(m \times n), \mB\;(n \times m) : \tr(\mA\mB)= \tr(\mB\mA)$\\
This holds more generally for any cyclical permutation of matrices $\mA_1,\ldots,\mA_n$ such that the product remains defined. 

\item\zlabel{it:20.7} $\mA,\mB\;(m \times m) : \det(\mA\mB) = \det(\mA)\det(\mB)$

\item\zlabel{it:20.8} $\mA\;(m \times n):\rk(\mA)=\rk(\mA^\top)=\rk(\mA^\top\mA)=\rk(\mA\mA^\top)$

\item\zlabel{it:20.9} $\mA\;(m\times n)$: $\mA^\top \mA$ and $\mA\mA^\top$ are symmetric positive semidefinite.\\
If, furthermore, $\rk(\mA)=n$, then $\mA^\top \mA$ is positive definite.

\end{enumerate}
\endprop
\end{frame}

\begin{frame}{Rules (3)}
\protect\hypertarget{rules-3}{}
\prop\zlabel{rules3}

For square \(\mA\; (m\times m)\) the following hold:

\begin{enumerate}
 \item\zlabel{it:21.1} $\tr(\mA^\top) = \tr(\mA)$
 
 \item\zlabel{it:21.2} $\det(\mA^\top) = \det(\mA)$ 
 
 \item\zlabel{it:21.3} $(\mA^\top)^{-1} = (\mA^{-1})^\top$, if $\mA$ is nonsingular. 

  \item\zlabel{it:21.4} $\det(\mA^{-1})= 1/\det(\mA)$ 
  
  \item\zlabel{it:21.5} $\lambda$ is eigenvalue of $\mA$ with eigenvector $\vv$ $\Rightarrow$ $1/\lambda$ is an eigenvalue of $\mA^{-1}$ with eigenvector $\vv$.

  \item\zlabel{it:21.6} $ \rk(\mA) = m \Leftrightarrow \mA^{-1}$ exists
    
  \item\zlabel{it:21.7} $\rk(\mA) =m \Rightarrow \mA^+ = \mA^{-1}$

\end{enumerate}
\endprop
\end{frame}

\begin{frame}{Orthogonal Matrices}
\protect\hypertarget{orthogonal-matrices}{}
\prop

\(\mA\) is orthogonal \(\Leftrightarrow\)

\begin{enumerate}
 \item\zlabel{it:22.1}  $\mA\mA^\top = \mI_m \Leftrightarrow \mA^\top \mA = \mI_m$. 
 
 \item\zlabel{it:22.2} $\mA^\top$ is orthogonal. 

\item\zlabel{it:22.3} $\mA^{-1}$ is orthogonal.
\end{enumerate}
\endprop
\end{frame}

\begin{frame}{Partitioned Inverses}
\protect\hypertarget{partitioned-inverses}{}
\prop\zlabel{partitionedinverses}

\(\mA\;(m \times m), \mB\;(m \times n), \mC\;(n \times m), \mD\;(n \times n)\)
: \[
\left(%
\begin{array}{cc}
  \mA & \mB \\
  \mC & \mD \\
\end{array}%
\right)^{-1}=\hspace{8cm}\]\[\left(%
\begin{array}{cc}
  \mA^{-1} + \mA^{-1}\mB(\mD-\mC\mA^{-1}\mB)^{-1}\mC\mA^{-1}
 & -\mA^{-1}\mB(\mD-\mC\mA^{-1}\mB)   \\
  -(\mD-\mC\mA^{-1}\mB)^{-1}\mC\mA^{-1}  & (\mD-\mC\mA^{-1}\mB)^{-1} \\
\end{array}%
\right),
\] provided the respective inverses exist. \endprop
\end{frame}

\begin{frame}{Partitioned Inverses}
\protect\hypertarget{partitioned-inverses-1}{}
\prop

For nonsingular \(\mA_i\;(m_i \times m_i)\), \(i = 1,\ldots, r\): \[
\left(%
\begin{array}{ccc}
  \mA_1 &  & \mZeros \\
   & \ddots &  \\
  \mZeros &  & \mA_r \\
\end{array}%
\right)^{-1}=\left(%
\begin{array}{ccc}
  \mA_1^{-1} &  &0 \\
   & \ddots &  \\
  0 &  & \mA_r^{-1} \\
\end{array}%
\right)
\] \endprop
\end{frame}

\begin{frame}{Relationship Between the Properties}
\protect\hypertarget{relationship-between-the-properties}{}
\prop
\begin{enumerate}
\item\zlabel{it:25.1} $\mB$ is nonsingular $\Rightarrow\ \det(\mB\mA\mB^{-1}) = \det(\mA)$

\item\zlabel{it:25.2} $\det(\mI_m) = 1$

\item\zlabel{it:25.3} $\mA\;(m \times m)$ with eigenvalues $\lambda_1,\ldots,\lambda_m : \det(\mA) =\prod_{j=1}^m
\lambda_j$

\item\zlabel{it:25.4} $\mA\;(m \times m) : \rk(\mA) < m\ \Leftrightarrow\ \det(\mA) = 0$

\item\zlabel{it:25.5} $\rk(\mA)=m\ \Leftrightarrow\ \det(\mA) \neq 0$

\item\zlabel{it:25.6} $\mA$ is singular $\Leftrightarrow\ \det(\mA) = 0$

\item\zlabel{it:25.7} If $\mA$ and $\mD$ are square, then 
\[
\left|%
\begin{array}{cc}
  \mA & \mB \\
  0 & \mD \\
\end{array}%
\right|=|\mA||\mD|
\] 

\item\zlabel{it:25.8} If $\mA$ and $\mD$ are square and $\mA^{-1}$ exists, then 
\[
\left|%
\begin{array}{cc}
  \mA & \mB \\
  \mC & \mD \\
\end{array}%
\right|=|\mA||\mD-\mC\mA^{-1}\mB|
\]
\end{enumerate}
\endprop
\end{frame}

\begin{frame}{Relationship Between the Properties}
\protect\hypertarget{relationship-between-the-properties-1}{}
\prop\label{propsidempotent}\zlabel{propsidempotent} For idempotent
\(\mA \;(m \times m)\):

\begin{enumerate} 
\item\zlabel{it:26.1} $\rk(\mA) = m\ \Rightarrow\ \mA = \mI_m$

\item\zlabel{it:26.2} $\rk(\mA) = \tr(\mA)$

\item\zlabel{it:26.3} $\rk(\mI_m-\mA) = m-\rk(\mA)$
\end{enumerate}
\endprop
\end{frame}

\begin{frame}{Relationship Between the Properties}
\protect\hypertarget{relationship-between-the-properties-2}{}
\prop\zlabel{relprop3}

For square \(\mA\;(m \times m)\):

\begin{enumerate}
 \item\zlabel{it:27.1} $\mA = \diag(a_{11},\ldots,a_{mm}) \Rightarrow
a_{11},\ldots,a_{mm}$ are the eigenvalues of $\mA$.\\
In fact, triangularity of $\mA$ suffices. 

\item\zlabel{it:27.2} $\mA$ is symmetric $\Rightarrow$ all eigenvalues of $\mA$ are real numbers. 

\item\zlabel{it:27.3} $\mA$ is idempotent $\Rightarrow$ all eigenvalues of $\mA$ are 0 or 1.

\item\zlabel{it:27.4} $\mA$ is singular $\Leftrightarrow$ 0 is an eigenvalue of $\mA$.

\item\zlabel{it:27.5} $\mA$ is nonsingular $\Leftrightarrow$ all eigenvalues of $\mA$ are nonzero.
 
\item\zlabel{it:27.6} $\mA$ symmetric: $\mA$ is positive definite $\Leftrightarrow$ all eigenvalues of $\mA$
are real and greater than 0. 

\item\zlabel{it:27.7} $\mA$ symmetric: $\mA$ is positive semidefinite $\Leftrightarrow$ all eigenvalues of $\mA$ are real and greater than or equal to 0. 

\item\zlabel{it:27.8} $\mA$ has eigenvalues $\lambda_1,\ldots,\lambda_m : \tr(\mA) = \sum_{j=1}^m\lambda_j$.\note{Proof goes by Schur decomposition of a square matrix and flipping the trace, see L\"utkepohl, 5.2.3(2)}
\end{enumerate}
\endprop
\end{frame}

\begin{frame}{Relationship Between the Properties}
\protect\hypertarget{relationship-between-the-properties-3}{}
\prop\zlabel{prop:pd and det}

\(\mA\) positive (semi) definite \(\Leftrightarrow\) \[ 
\det\left(\begin{array}{ccc}
  a_{11} & \cdots & a_{1i} \\
  \vdots & \ddots & \vdots \\
  a_{i1} & \cdots & a_{ii}
\end{array}%
\right)>0\;(\geqslant 0)\;\;\textrm{ for } i = 1,\ldots,m.
\] \(\mA\) negative (semi) definite \(\Leftrightarrow\) \[
\det\left(\begin{array}{ccc}
  a_{11} & \cdots & a_{1i} \\
  \vdots & \ddots & \vdots \\
  a_{i1} & \cdots & a_{ii}
\end{array}%
\right)\begin{cases}>0\;(\geqslant 0) &i \textrm{
even}\\<0\;(\leqslant 0) &i \textrm{ odd}\end{cases}\;\;\textrm{
for } i = 1,\ldots,m.
\] \endprop
\end{frame}

\begin{frame}{Spectral Decomposition of a Symmetric Matrix}
\protect\hypertarget{spectral-decomposition-of-a-symmetric-matrix}{}
\thmn\zlabel{spectraldecomp}

Let \(\mA\;(m \times m)\) be symmetric with eigenvalues
\(\lambda_1,\ldots,\lambda_m\). Then, \[ 
  \mA= \mU\mLambda \mU^\top,
\] where \(\mLambda = \diag(\lambda_1,\ldots,\lambda_m)\) and \(\mU\) is
the real orthogonal \((m \times m)\) matrix whose columns are the
orthonormal eigenvectors \(\vv_1,\ldots,\vv_m\) of \(\mA\)
(i.e.~\(\vv_i^\top\vv_j=1\) if \(i=j\) and \(0\) else) associated with
\(\lambda_1,\ldots,\lambda_m\). In other ``words,'\,' \[
  \mA=\sum_{j=1}^m\lambda_j\vv_j\vv_j^\top.
\] \endthmn
\end{frame}

\begin{frame}{Cholesky Decomposition}
\protect\hypertarget{cholesky-decomposition}{}
\thmn\label{cholesky}\zlabel{cholesky} \(\mA\;(m \times m)\) symmetric
positive definite:

There exists a unique lower (upper) triangular \((m \times m)\) matrix
\(\mB\) with real positive diagonal elements such that
\(\mA = \mB\mB^\top\). \endthmn

\xmpl \(\mA\;(2\times 2)\)

\begin{eqnarray*}
\left(%
\begin{array}{cc}
  a_{11} & a_{12} \\
  a_{12} & a_{22} \\
\end{array}%
\right)&=&\left(%
\begin{array}{cc}
  b_{11} & 0 \\
  b_{12} & b_{22} \\
\end{array}%
\right)\left(%
\begin{array}{cc}
  b_{11} & b_{12} \\
  0 & b_{22} \\
\end{array}%
\right) \Rightarrow \\[.5cm]
b_{11}&=&\sqrt{a_{11}},\\b_{12}&=&a_{12}/b_{11},\\b_{22}&=&\sqrt{a_{22}-b_{12}^2}.
\end{eqnarray*} \endxmpl
\end{frame}

\begin{frame}{Matrix Calculus}
\protect\hypertarget{matrix-calculus}{}
Let \(f(\vx)\) be a differentiable real-valued function of the
\((m \times 1)\) vector \(\vx = (x_1,\ldots,x_m)^\top\).

\defn[Gradient]\zlabel{def:Gradient}

The \hil{gradient} of \(f(\cdot)\) is \[
  \frac{\partial f}{\partial \vx}:=\left(%
\begin{array}{c}
  \frac{\partial f}{\partial x_1} \\
  \vdots \\
  \frac{\partial f}{\partial x_m} \\
\end{array}%
\right)\;\;(m\times 1)
\] and \[
  \frac{\partial f}{\partial \vx^\top}:=\left(
  \frac{\partial f}{\partial x_1} , \cdots , \frac{\partial f}{\partial x_m}\right)\;\;(1\times m)
\] \enddefn
\end{frame}

\begin{frame}{Matrix Calculus}
\protect\hypertarget{matrix-calculus-1}{}
\defn[Hessian]The \hil{Hessian} of \(f(\cdot)\) is \[
  \frac{\partial^2 f}{\partial \vx\partial \vx^\top}:=\left(%
\begin{array}{ccc}
  \frac{\partial^2 f}{\partial x_1\partial x_1}&\cdots&\frac{\partial^2 f}{\partial x_1\partial x_m} \\
  \vdots &\ddots&\vdots\\
  \frac{\partial^2 f}{\partial x_m\partial x_1}&\cdots&\frac{\partial^2 f}{\partial x_m\partial x_m}\\
\end{array}%
\right)\;\;(m\times m)
\] \enddefn
\end{frame}

\begin{frame}{Matrix Calculus}
\protect\hypertarget{matrix-calculus-2}{}
\defn[Jacobian]Let \(\vy(\vx)=(y_1(\vx),\ldots, y_n(\vx))^\top\) be an
\((n \times 1)\) vector of differentiable functions of the
\((m \times 1)\) vector \(\vx= (x_1,\ldots, x_m)^\top\). That is,
\(\vy\) is a function mapping a subset of \(\mathbb{R}^m\) on a subset
of \(\mathbb{R}^n\).\medskip

The \hil{Jacobian} of \(\vy(\cdot)\) is defined as \[
  \frac{\partial \vy}{\partial \vx^\top}:=\left(%
\begin{array}{ccc}
  \frac{\partial y_1}{\partial x_1}&\cdots&\frac{\partial y_1}{\partial x_m} \\
  \vdots &\ddots&\vdots\\
  \frac{\partial y_n}{\partial x_1}&\cdots&\frac{\partial y_n}{\partial x_m}\\
\end{array}%
\right)\;\;(n\times m)
\] \enddefn
\end{frame}

\begin{frame}{Matrix Calculus}
\protect\hypertarget{matrix-calculus-3}{}
\prop\zlabel{matrdiff}

We have

\begin{enumerate}
\item $\vx,\;\va\;(m \times 1)$ \[\frac{\partial \va^\top \vx}{\partial \vx}=\va\] 

\item  $\vx\; (m \times 1),\;\mA\;(m \times m)$
\[
  \frac{\partial \vx^\top \mA \vx}{\partial \vx}=(\mA+\mA^\top)\vx
\] 
(Note how the formula simplifies if $\mA$ is symmetric.)

\item
\[
  \frac{\partial^2 \vx^\top \mA \vx}{\partial \vx\partial \vx^\top}=(\mA+\mA^\top)
\]
\end{enumerate}
\endprop
\end{frame}




\end{document}